% =====================================================================
% Paper 1 v64 (research-article rewrite of paper1_assessment_separation_v62/v63.tex;
% repo v63 = v62 + one paragraph on the Lean layer, incorporated in Section 4.7)
% Every weighting licenses a plan, no plan satisfies the floors:
% a quantifier-order separation between compensatory and noncompensatory
% transition assessment
% Author: Amin Abaee (Independent Researcher, Tehran, Iran)
% Target: Journal of Multi-Criteria Decision Analysis (special issue
% "Better Decisions for a Better Tomorrow")
%
% v64 (2026-10-14): restructured as a research article. Results renumbered in
% presentation order (map in CHANGELOG_v62_to_v64.md). The finite-menu geometric
% theorem now leads the results; the witness datum, the blend/time-sharing pair,
% the rescue threshold, the substitutability ladder and the common-shock verdict
% follow as instances. New Figures 1 (finite-menu geometry) and 4 (ladder).
% Terminology guide, practitioner checklist, minimum report, data requirements,
% the full substitutability-spectrum and common-shock extensions, the complete
% benchmark, the check-list enumeration and the framework extensions are in the
% Supplementary Information. No theorem, number or proof was changed.
% Compiles with pdfLaTeX / XeLaTeX / LuaLaTeX / tectonic (elsarticle class).
% =====================================================================
\documentclass[preprint,11pt,authoryear]{elsarticle}

\usepackage{amsmath,amssymb,amsthm}
\usepackage{booktabs}
\usepackage{array}
\usepackage{tabularx}
\usepackage{graphicx}
\graphicspath{{./}{figs/}{paper1_assessment_separation_v64_src/figs/}}
\usepackage{enumitem}
\usepackage[colorlinks=true,urlcolor=blue,linkcolor=blue,citecolor=blue]{hyperref}
\emergencystretch=3em
\setlength{\tolerance}{2000}

\newtheorem{theorem}{Theorem}
\newtheorem{proposition}{Proposition}
\newtheorem{lemma}{Lemma}
\newtheorem{corollary}{Corollary}
\theoremstyle{definition}
\newtheorem{definition}{Definition}
\theoremstyle{remark}
\newtheorem{remark}{Remark}

\newcommand{\Etyp}{E_{\mathrm{typ}}}
\newcommand{\Eend}{E_{\mathrm{end}}}
\newcommand{\Ephys}{E_{\mathrm{tube,phys}}}
\newcommand{\Vtyp}{\mathcal{V}_{\mathrm{typ}}}
\newcommand{\Vweak}{\mathcal{V}_{\mathrm{weak}}}
\newcommand{\Vphys}{\mathcal{V}_{\mathrm{phys}}}
\newcommand{\Vend}{\mathcal{V}_{\mathrm{end}}}
\newcommand{\Wp}{W_+}
\newcommand{\FP}{\mathrm{FP}_{\mathrm{agg}}}
\newcommand{\Tube}{\mathrm{Tube}}
\newcommand{\Succ}{\mathrm{Succ}}
\newcommand{\End}{\mathrm{End}}
\newcommand{\conv}{\operatorname{conv}}
\newcommand{\Rp}{\mathbb{R}^n_+}
\newcommand{\Q}{\mathcal{Q}}
\newcommand{\act}[1]{\textsc{#1}}

\journal{Journal of Multi-Criteria Decision Analysis}

\begin{document}

\begin{frontmatter}

\title{Every weighting licenses a plan, no plan satisfies the floors: a quantifier-order separation between compensatory and noncompensatory transition assessment}

\author[1]{Amin Abaee\corref{cor1}}
\ead{amin_abaee@ut.ac.ir}
\cortext[cor1]{Corresponding author. ORCID: 0000-0002-0019-1842.}
\affiliation[1]{organization={Independent Researcher}, city={Tehran}, country={Iran}}

\begin{abstract}
Composite indices sort transitions into accepted and rejected under one of two acceptance semantics. Under weight-adaptive (compensatory) acceptance a transition is accepted if every weighting of the index licenses some plan, possibly a different plan for each weighting; under common-plan (noncompensatory) acceptance one plan must keep every criterion above its floor along the whole path. We compare the two as robust finite-horizon certification problems on a common datum with path-wise constraints and worst-case disturbances. At any fixed trajectory the full-cone aggregate loses nothing, yet the semantics separate, because ``some plan for every weighting'' and ``one plan for all weightings'' do not commute. For any finite menu of plans governed by required-margin vectors, the acceptance gap is the part of the convex hull of the margins that no single plan dominates, so it vanishes exactly when the menu is convexified. On an explicit rational datum the gap has nonempty interior; its states split, at a closed-form reserve threshold, into those a staged plan can rescue and those no available plan completes safely, with exact weight-ratio thresholds giving the preference range over which each plan is licensed. Blending plans closes the gap; alternating them in time does not. Every positive-elasticity CES index falsely certifies an exact interval of gap states, with critical floors $1<5/4<\sqrt2<\varphi<\sqrt3<\dots\to2$, and a shared shock reduces the gap by an exact cut without removing it. A fishery transition anchored to a regulated biomass limit realises the datum. Where separately binding floors matter, the aggregate alone cannot tell the rescuable from the impossible: per-floor reporting along the path, or a verified common plan, is a detection requirement.
\end{abstract}

\begin{keyword}
multi-criteria decision analysis \sep compensatory and noncompensatory aggregation \sep composite indicators \sep scalarization \sep robustness \sep sustainability assessment \sep viability
\end{keyword}

\end{frontmatter}

\section{Introduction}\label{sec:intro}

\subsection{Compensability as an acceptance question}

Multi-criteria assessment of a sustainability transition faces a choice that the decision-analysis literature has studied since its foundations: whether the criteria may compensate one another. Multi-attribute value theory builds substitution rates between criteria into its trade-off structure \citep{KeeneyRaiffa1976}; the outranking tradition constrains compensation through concordance and veto \citep{Roy1996,Vincke1992}; and the decision-aiding treatment of sustainable development centres on incommensurable criteria \citep{Munda2005,MartinezAlier1998}. In environmental assessment the same choice appears as the distinction between weak sustainability, in which capital forms are substitutable through prices and an aggregate must not decline \citep{WorldBank2011,Neumayer2013,Boos2015}, and strong sustainability, in which certain stocks are separately binding \citep{Daly1990,Ekins2003}; the indicator theory of the two positions is well developed \citep{Hanley1999,Martinet2011,CairnsMartinet2014,DoyenGajardo2020,Usubiaga2025}, as is the compensability analysis of multi-criteria methods used for sustainability assessment \citep{Cinelli2014,Schar2025}.

Nearly all of this work classifies aggregation behaviour at a fixed profile of criterion values. This paper asks a different question, about \emph{transitions}: when an authority must carry out a plan over time, under disturbances, with several criteria required to hold along the entire path, when can a compensatory aggregate certify a transition that a noncompensatory assessment rejects? The question is not about the choice of weights, whose methodology is itself a developed field \citep{Nardo2008,Becker2017}, and it is not about the known static limitation of weighted sums, which cannot reach nonconvex parts of a Pareto front \citep{DasDennis1997}. It is about the \emph{acceptance semantics} applied to a composite reading of a transition: the rule by which the reading sorts the transition into the accepted or rejected category, the sorting problematic of decision aiding \citep{Roy1996}.

\subsection{Two acceptance semantics}

Fix a transition datum: a state carrying typed floors $s_i\ge0$ (criteria normalised so that zero is the binding level), a finite menu of plans, a disturbance set, and \emph{exact-tube} semantics under which a plan is safe only if every state visited along it, not merely its endpoint, satisfies the constraints. For each nonnegative weighting $w$ in the full cone $\Wp=\Rp\setminus\{0\}$ let $E_w(z)$ be the set of plans admissible at state $z$ under the aggregate floor $w\cdot s\ge0$, and let $\Etyp(z)$ be the set admissible under the typed floors taken separately. Two acceptance criteria differ by quantifier order:
\begin{itemize}[leftmargin=2em]
\item \emph{common-plan acceptance} (noncompensatory, coordinate-wise): $\exists a\ \forall w:\ a\in E_w(z)$;
\item \emph{weight-adaptive acceptance} (compensatory, scalarized): $\forall w\ \exists a_w:\ a_w\in E_w(z)$.
\end{itemize}
The first implies the second; the converse fails. Table~\ref{tab:protocols} places the two among the four protocols an index can serve. In the vocabulary of decision making under uncertainty they differ in \emph{when the plan is fixed relative to the weight}: common-plan acceptance commits before the weighting is known, weight-adaptive acceptance lets the plan adapt once it is observed, the here-and-now versus wait-and-see distinction of adjustable robust optimisation \citep{BenTal2004} with the assessment weight in the role of the uncertain parameter. The separation studied here is Protocol~2 against Protocol~4; it is a statement about a certification protocol applied to a composite index, not an objection to composite indices as such.

\begin{table}[!htbp]
\centering\small
\caption{Four acceptance protocols for a composite index on a transition datum. The separation studied in this paper is Protocol 2 against Protocol 4; Protocol 3 coincides with Protocol 4 on the full cone (Proposition~\ref{prop:hierarchy}(ii)).}
\label{tab:protocols}
\begin{tabularx}{\textwidth}{>{\raggedright\arraybackslash}p{0.25\textwidth} >{\raggedright\arraybackslash}p{0.27\textwidth} X}
\toprule
Protocol & Quantifier form & Reading \\
\midrule
1. Fixed-weight assessment & $a\in E_w(z)$ at a declared $w$ & one index, one selected plan: ordinary fixed-weight certification \\
2. Weight-adaptive (per-weight) acceptance & $\forall w\ \exists a_w\in E_w(z)$ & every weighting can be paired with some plan, possibly weighting-dependent: the compensatory reading, $\Vweak=\bigcap_w\mathcal V_w$ \\
3. Common-plan robustness & $\exists a\ \forall w:\ a\in E_w(z)$ & one plan serves every weighting \\
4. Coordinate-wise (typed) acceptance & $\exists a:\ a\in\Etyp(z)$ & separately binding floors enforced path-wise: the noncompensatory reading \\
\bottomrule
\end{tabularx}
\end{table}

A concrete form of the question: a public authority must carry out a restructuring while satisfying two minimum-service constraints. Rapid restructuring exposes one constraint to a temporary drawdown; a gradual plan shifts the burden to the other; a staged plan avoids both but requires a reserve to finance it. An evaluator who aggregates the two constraints into one weighted index may approve the transition for \emph{every} relative price, each price licensing a different plan, while an evaluator who requires each constraint to hold along the whole path approves none of the available plans. The results below show that this divergence is generic in a precise geometric sense, that it is bounded by an exact reserve threshold, that it is closed by blending plans but not by alternating them, and that it survives the replacement of the linear index by any positive-elasticity aggregator.

\subsection{Contributions}

The paper makes five contributions, stated in the order of the text.
\begin{enumerate}[leftmargin=2em,label=(\roman*)]
\item \emph{Localisation and hierarchy} (Section~\ref{sec:general}). On the full cone the typed operator is exactly the common-plan intersection, $\Etyp(z)=\bigcap_{w\in\Wp}E_w(z)$, and at any fixed trajectory the aggregate is lossless (Lemma~\ref{lem:cone}); the chain typed $\subseteq$ scalarized $\subseteq$ physical $\subseteq$ endpoint holds for every weight, and restricting the weight family can only enlarge the compensatory accepted set (Propositions~\ref{prop:hierarchy}--\ref{prop:monotone}). The separation is therefore dynamic -- a property of the quantifier over plans -- and the full cone is its strictest case.
\item \emph{A geometric characterisation for any finite menu} (Theorem~\ref{thm:menu}). For plans whose coordinate-wise admissibility is governed by a required-margin vector $d_a$, weight-adaptive acceptance is $\conv(\mathcal D)+\Rp$ and common-plan acceptance is $\bigcup_a(d_a+\Rp)$; the acceptance gap is the part of the convex hull that no single plan dominates, and convexifying the menu closes it.
\item \emph{An explicit rational witness} (Theorem~\ref{thm:witness}). On a four-plan datum the gap is a region with nonempty interior, equal to the \emph{impossibility region} $I$; the aggregate-versus-floor discrepancy region splits at the reserve level $x=1$ into $I$ and a \emph{rescue set} $R$ that a reserve-financed staged plan already serves; exact weight-ratio thresholds $\rho_1=(2-s_1)/s_2$ and $\rho_2=s_1/(2-s_2)$ give the preference range over which each plan is licensed; and the minimal reserve increment that converts impossibility into rescue is $\kappa^*(z)=(1-x)\,\mathbf 1[z\notin\Vtyp]$ (Proposition~\ref{prop:rescue}). The separation persists under hold-prefix extension of the horizon and can be erased by a later stage with coincident terminal sets (Proposition~\ref{prop:persist}).
\item \emph{What closes the gap and what does not} (Theorem~\ref{thm:blend}, Proposition~\ref{prop:timeshare}). Fractional blending of the primitive plans -- menu convexification at the level of control flows -- makes the typed accepted set coincide exactly with the compensatory region, so the gap is a property of the finite deterministic menu, not of either doctrine; discrete time-sharing closes it only where both dips are separately survivable.
\item \emph{Robustness of the separation to the aggregator and to the disturbance class} (Theorem~\ref{thm:ladder}, Proposition~\ref{prop:commonshock}). Replacing the linear index by a constant-elasticity-of-substitution (CES) aggregator with elasticity $\sigma$ changes the gap but does not remove it: the accepted rungs are nested, each positive-elasticity member falsely certifies an exact interval of gap states with critical floors $1<5/4<\sqrt2<\varphi<\sqrt3<\dots\to2$, and only the Leontief member ($\sigma=0$), which is the typed operator, is uniformly safe. Under a shared environmental shock striking every plan the gap survives, reduced by an exact $1/2$-cut, with the licensing thresholds unshifted on the surviving region; the reserve route then fails on an exact band of states.
\end{enumerate}
Every continuum statement is proved (\ref{app:general}--\ref{app:witness}; Supplementary Information, SI~S2) and every quoted number is an exact rational re-derived by deposited verification code; the core theorems are additionally formalised in Lean~4 (Section~\ref{sec:verification}). A resource-transition benchmark (Section~\ref{sec:benchmark}) realises the datum as a fishery transition with units anchored to a regulated biomass limit.

\paragraph{What is not claimed} No ranking of weak and strong sustainability as doctrines; no proof that any particular set of boundaries ought to be treated noncompensatorily; no empirical result about any resource system; no infinite-horizon, stochastic or partially observed extension. The impossibility region is a negative certificate relative to the specified plan menu, and the separation is exhibited for a declared disturbance class.

\subsection{Position in the literature}

The backward constructions used here are typed instances of robust-predecessor, capture-basin and reach--avoid constructions \citep{Aubin1991,Aubin2011,SaintPierre1994,Lygeros1999,AubinCatte2002}; the hierarchy of Proposition~\ref{prop:hierarchy}(i) is constraint-set monotonicity of viability kernels \citep{Aubin1991,Frankowska1989}; viability methods are established in natural-resource management \citep{DeLaraDoyen2008,Doyen2012}. The blend-collapse pair is the pure-versus-mixed-strategy structure of the assessor--planner game \citep{vonNeumann1928,Sion1958}, whose reach--avoid form is the discriminating kernel of \citet{Cardaliaguet1996} \citep[numerically,][]{Cardaliaguet1999}. In composite-indicator practice the critique of weights as value judgements presented as measurement \citep{Hickel2020,Fischer2022} and the construction of indices without weights from Pareto fronts \citep{Gao2023} both concern the index; the result here is orthogonal -- even with weights elicited faithfully, the quantifier order produces the gap. The exactness statement of Lemma~\ref{lem:cone} is, in the vocabulary of the companion obstruction calculus \citep{Abaee2026obs}, the statement that the typed property is constant on the fibres of the full-cone observation at a fixed trajectory; Theorem~\ref{thm:witness} shows that this exactness does not survive the quantifier over plans. We do not assert priority over the elementary fact $\exists a\,\forall w\ne\forall w\,\exists a_w$; to our knowledge no explicit exact-tube witness of the dynamic separation in a compensatory/noncompensatory assessment setting has been reported, but whether an equivalent separation appears in adjacent literatures is not established either way.

